% !TeX program = xelatex
\documentclass[8pt]{beamer}

\usepackage[english]{minimalistslides}
\usepackage{booktabs}

% This starter keeps upstream's own default logo (figs/logo.jpg, Southeast
% University's crest) as the worked example of \setthemelogo -- the
% template is explicitly designed to be rebranded (see upstream's README:
% "Fast institutional rebranding through logo, metadata, and color
% overrides"), so the institution name below matches the bundled crest
% rather than inventing an unrelated one.
\setthemename{Southeast University}
\setthemelogo{figs/logo.jpg}

\renewcommand{\presentationtitle}{Battery Degradation Under Fast Charging}
\renewcommand{\presentationauthor}{Sofia Mendes}
\renewcommand{\presentationinstitution}{\themename}
\renewcommand{\presentationdate}{\today}

\begin{document}

\makeboxedcover

\begin{frame}{Outline}
  \tableofcontents
\end{frame}

\section{Motivation}

\makesection

\begin{frame}{Why Fast Charging Matters}
  Fast charging cuts wait time from hours to minutes, but every extra
  \highlight{amp} pushed through a cell accelerates capacity loss.

  \begin{itemize}
    \item Consumer demand for sub-20-minute charging keeps growing
    \item Degradation mechanisms are well studied at low rates, less so above 3C
    \item This work isolates the fast-charging contribution specifically
  \end{itemize}
\end{frame}

\section{Method}

\makesection

\begin{frame}{Test Protocol}
  \begin{table}
    \small
    \centering
    \begin{tabular}{lrr}
      \toprule
      Group & Charge rate & Cycles \\
      \midrule
      Baseline & 1C & 500 \\
      Fast & 3C & 500 \\
      Ultra-fast & 5C & 500 \\
      \bottomrule
    \end{tabular}
  \end{table}
  Capacity measured every 25 cycles under identical discharge conditions.
\end{frame}

\begin{frame}{What We Tracked}
  \begin{itemize}
    \item Capacity fade relative to first-cycle baseline
    \item Internal resistance growth
    \item Post-mortem electrode imaging on a subset of cells
  \end{itemize}
\end{frame}

\section{Results}

\makesection

\begin{frame}{Capacity Fade by Charge Rate}
  \begin{table}
    \small
    \centering
    \begin{tabular}{lrrr}
      \toprule
      Rate & 100 cyc. & 300 cyc. & 500 cyc. \\
      \midrule
      1C & 98.4\% & 94.1\% & 89.7\% \\
      3C & 96.2\% & 88.3\% & 79.5\% \\
      5C & \stronghighlight{93.8\%} & \stronghighlight{79.6\%} & \stronghighlight{64.2\%} \\
      \bottomrule
    \end{tabular}
  \end{table}
  \stronghighlight{5C charging} loses capacity roughly three times faster
  than the 1C baseline by cycle 500.
\end{frame}

\begin{frame}{Where the Loss Comes From}
  Post-mortem imaging on the 5C group shows visible lithium plating on the
  anode edge, concentrated near the tab -- consistent with localized
  current density rather than uniform aging.
\end{frame}

\section{Conclusion}

\makesection

\begin{frame}{Takeaways}
  \begin{itemize}
    \item Fast charging above 3C accelerates fade disproportionately
    \item Plating, not uniform SEI growth, drives most of the extra loss
    \item Edge-focused current density suggests a tab-design fix is worth testing
  \end{itemize}
\end{frame}

\begin{frame}{Thank You}
  \centering
  \Large Questions?\\[1em]
  \normalsize sofia.mendes@example.edu
\end{frame}

\end{document}
